% Options for packages loaded elsewhere
\PassOptionsToPackage{unicode}{hyperref}
\PassOptionsToPackage{hyphens}{url}
\documentclass[
]{article}
\usepackage{xcolor}
\usepackage{amsmath,amssymb}
\setcounter{secnumdepth}{-\maxdimen} % remove section numbering
\usepackage{iftex}
\ifPDFTeX
  \usepackage[T1]{fontenc}
  \usepackage[utf8]{inputenc}
  \usepackage{textcomp} % provide euro and other symbols
\else % if luatex or xetex
  \usepackage{unicode-math} % this also loads fontspec
  \defaultfontfeatures{Scale=MatchLowercase}
  \defaultfontfeatures[\rmfamily]{Ligatures=TeX,Scale=1}
\fi
\usepackage{lmodern}
\ifPDFTeX\else
  % xetex/luatex font selection
\fi
% Use upquote if available, for straight quotes in verbatim environments
\IfFileExists{upquote.sty}{\usepackage{upquote}}{}
\IfFileExists{microtype.sty}{% use microtype if available
  \usepackage[]{microtype}
  \UseMicrotypeSet[protrusion]{basicmath} % disable protrusion for tt fonts
}{}
\makeatletter
\@ifundefined{KOMAClassName}{% if non-KOMA class
  \IfFileExists{parskip.sty}{%
    \usepackage{parskip}
  }{% else
    \setlength{\parindent}{0pt}
    \setlength{\parskip}{6pt plus 2pt minus 1pt}}
}{% if KOMA class
  \KOMAoptions{parskip=half}}
\makeatother
\setlength{\emergencystretch}{3em} % prevent overfull lines
\providecommand{\tightlist}{%
  \setlength{\itemsep}{0pt}\setlength{\parskip}{0pt}}
\usepackage{bookmark}
\IfFileExists{xurl.sty}{\usepackage{xurl}}{} % add URL line breaks if available
\urlstyle{same}
% fallback for those not using the hyperref driver hyperxmp:
\makeatletter
\@ifundefined{xmpquote}{\newcommand{\xmpquote}[1]{#1}}{}
\makeatother
\hypersetup{
  hidelinks,
  pdfcreator={LaTeX via pandoc}}

\author{}
\date{}

\begin{document}

{
\setcounter{tocdepth}{3}
\tableofcontents
}
\section{Source Provenance Policy}\label{source-provenance-policy}

\subsection{1. Purpose and scope}\label{purpose-and-scope}

This policy protects Password Bonobo's original implementation work and
records how compatibility research becomes reviewable, neutral project
knowledge. It applies to Bonobo-authored source, tests, fixtures,
documentation, dependencies, assets, and research records. It does not
grant permission to reuse material owned by another party.

\subsection{2. Allowed authoritative
inputs}\label{allowed-authoritative-inputs}

Product implementation and tests may use these authoritative inputs:

\begin{itemize}
\tightlist
\item
  Approved Bonobo specifications and decisions.
\item
  Official PasswordSafe format documentation and other official format
  documentation applicable to the work.
\item
  The approved neutral behavior dossier, which describes observable
  behavior without reproducing implementation text.
\item
  Bonobo synthetic tests and fixtures that passed the intake process in
  this policy.
\end{itemize}

These inputs are the handoff boundary for implementation. A developer
must not translate external implementation material into Bonobo code,
tests, comments, identifiers, or document structure.

\subsection{3. External research boundary and read-only checkout
rules}\label{external-research-boundary-and-read-only-checkout-rules}

Password Gorilla may be studied only in an external, pinned, read-only
research checkout outside this repository. Gorilla source is evidence
about behavior, not implementation material for Bonobo. Researchers may
inspect the pinned checkout for neutral observations, but must not
modify it, add Bonobo annotations to it, copy it into this repository,
or include it in a product build.

Research findings enter Bonobo through cited, neutral observations in
the compatibility dossier. Product development uses the approved Bonobo
documents, official format documentation, and synthetic tests rather
than the external source.

\subsection{4. Prohibited copying
categories}\label{prohibited-copying-categories}

Do not copy or adapt external source, comments, identifiers, file
organization, control flow, UI assets, translations, test fixtures,
screenshots, or other copyrightable implementation expression. Do not
quote external source fragments in product code or compatibility
documents. Do not import upstream vault files unless they pass the
fixture intake process and receive explicit approval.

\subsection{5. Evidence citation format}\label{evidence-citation-format}

Each compatibility evidence record must identify, in neutral prose:

\begin{itemize}
\tightlist
\item
  Revision: the pinned upstream revision.
\item
  Location: repository-relative path and either a line range or a test
  name.
\item
  Evidence kind: for example, source inspection, test observation,
  documentation, or change history.
\item
  Neutral observation: the observable behavior or unresolved question,
  without copied source text.
\end{itemize}

Use the cited material only to support the observation. A citation is
not permission to reproduce or implement the source expression it
references.

\subsection{6. Fixture intake}\label{fixture-intake}

Before a fixture is accepted, record its origin and license, establish
that its data is synthetic or otherwise approved for redistribution,
scan it for sensitive data, and obtain the required approval. The intake
record must state the reviewer or approving authority only when that
role and decision have been formally defined; it must not invent one.
Fixtures containing real credentials, personal vault data, tokens, or
provider metadata are prohibited.

\subsection{7. Dependency and asset
ledger}\label{dependency-and-asset-ledger}

Every dependency and nontrivial asset must have a ledger entry before
use. The entry must identify its origin, version or immutable revision,
license or terms, intended use, distribution implications, provenance
evidence, and required review status. The ledger must flag dependencies
and assets that are incompatible with an iOS distribution decision or
with the project's no-copy boundary.

\subsection{8. Upstream revision
updates}\label{upstream-revision-updates}

Updating an external research revision requires a separate review, a
newly pinned immutable revision, a provenance entry, and a neutral
dossier delta. Keep the checkout external and read-only throughout the
update. Do not replace an existing observation merely because a later
upstream revision differs; record the scope of the revised evidence and
revisit affected compatibility claims.

\subsection{9. Suspected copied-expression incident
response}\label{suspected-copied-expression-incident-response}

If copied expression is suspected, stop distributing and modifying the
affected material, preserve the relevant provenance records, and
restrict discussion to information safe to share. Escalate the concern
to the repository owner for review. Do not resolve the incident by
rewriting history, deleting evidence, or publishing the suspected
material. Resume work only after the affected material, evidence, and
required corrective action have been reviewed.

\end{document}
